% =============================================================================
%  VII. DISCUSSION: MAIN FINDINGS AND LIMITATIONS
% =============================================================================
\begin{frame}{VII. Discussion: Main Findings and Limitations}
\begin{atucols}[t]
\begin{column}{0.50\textwidth}
\textbf{Main findings}
\begin{enumerate}\itemsep2pt\scriptsize
  \item \textbf{Treatment-dependent membership and onset timing:} the
        pre-registered onset rule made H3 non-evaluable; standard timing moved
        \pcFNMovedEarlier{} onsets earlier.
  \item \textbf{Definition sensitivity:} prevalence ranged from
        \pcPctOnsetsA\% to \pcPctOnsetsC\% and case agreement was low
        ($\kappa = \pcKappaVsBA$), but the AUROC range was only
        \pcSpreadLabel.
  \item \textbf{Limited clinical utility:} the best Variant B operating point
        required \pcUtilityOptBurdenGbtB{} false alerts per true alert
        (benchmark 1.4).
\end{enumerate}
\vspace{2pt}
\begin{alertblock}{}
  \scriptsize Improving discrimination alone cannot resolve limitations
  introduced by outcome membership, onset timing and alert burden.
\end{alertblock}
\end{column}
\begin{column}{0.48\textwidth}
\textbf{Limitations}
\begin{itemize}\itemsep2pt\scriptsize
  \item Current-hour onset detection, not advance prediction.
  \item External Variant B analysis: \pcExtMicroCoverage\% culture coverage,
        \pcExtNEventsB{} events.
  \item Single-centre development; first ICU stays of at least 4\,h.
  \item The pre-registered onset rule differs from the standard rule;
        Variant F restores standard timing post hoc.
  \item Variant D is an unvalidated label whose components overlap the
        model inputs.
  \item Post-hoc analyses are labelled; the literature review is narrative.
\end{itemize}
\end{column}
\end{atucols}
\end{frame}


% =============================================================================
%  VIII. CONCLUSION: CONTRIBUTIONS AND FUTURE WORK
%  C-numbers match the contribution list in the thesis conclusion chapter.
% =============================================================================
\begin{frame}{VIII. Conclusion: Contributions and Future Work}
\begin{center}\renewcommand{\arraystretch}{1.1}\footnotesize
\begin{tabular}{cp{11.6cm}}
\toprule
C1 & Model class produced the largest AUROC range; label definition changed
     AUROC by only \pcSpreadLabel{} despite low case agreement
     ($\kappa = \pcKappaVsBA$). \\
C2 & The pre-registered onset rule leaves the pre-treatment window
     structurally empty; standard timing places \pcPreTreatOnsetsVarF{} test-set onsets
     there. \\
C3 & Maximum Variant B utility was \pcUtilityMaxGbtB, requiring
     \pcUtilityOptBurdenGbtB{} false alerts per true alert. \\
C4 & \pcHTwoNRobust{} covariates remained unstable after accounting for
     estimation uncertainty, all test-order indicators. \\
\bottomrule
\end{tabular}
\end{center}
{\scriptsize\color{MidGrey} Also: C5 significant label-sensitivity
contrasts involve missingness categories; C6 pre-registered hypotheses,
disclosed amendments, a fixed seed and values written mechanically from result
tables; C7 definitional versus empirical label
variation ($\kappa_{D,B}$, $\phi_{\text{anchor}}$).\par}

\vspace{4pt}
{\scriptsize\textbf{Future work:} labels less dependent on treatment
timestamps and adjudicated independently of the model inputs; lagged
predictors; SEP-1 and CDC Adult Sepsis Event definitions; prospective
multi-centre validation; utility-optimised training.\par}

\vspace{2pt}
\begin{block}{}
  \centering\footnotesize The main effect of label choice was on cohort
  membership and onset timing rather than on AUROC. Model choice affected
  AUROC more, and utility stayed close to the no-alert strategy.
\end{block}
\end{frame}
